\documentclass[aps,prb,floatfix,amsmath,amssymb,groupedaddress]{revtex4}
\usepackage{amsmath}
\usepackage{amsthm}
\newtheorem{prop}{Proposition}

\begin{document}
\providecommand{\half}{{\frac{1}{2}}}

\title{Extensions to ``Intro to Optical Lattices'': higher dimensions, optimal shuttles,\\
shaped illumination, and quasicrystals}
\author{M.~German (with AI assistance) --- a companion note to the public
manuscript and code of H.~Swan (github.com/orswan); independent work, not
endorsed by the original author}
\maketitle

This note extends the shuttling/Landau-Zener analysis of Swan's manuscript in
four directions. The dynamical claims (Secs.~I--III) were checked against
momentum-space Schr\"odinger integration at finite cutoff
(\texttt{lattice3D.py}, a $d$-dimensional generalization of his
\texttt{lattice1D.py}; natural units $\hbar = 1$, $2m = 1$ throughout), with
endpoint- and resolution-convergence spot checks; Secs.~IV--V rest on the
separate numerical models described there.

\section{The projection law in $d$ dimensions}

Consider a $d$-dimensional lattice $V(x) = \sum_c A_c \cos(b_c \cdot x)$ swept
by a frame trajectory $Q(t)$ with velocity $\dot Q$. Zone-edge degeneracies of
the 1D problem become \emph{Bragg planes}: the state $|q\rangle$ crosses
$|q + b\rangle$ where $|q - Q|^2 = |q + b - Q|^2$. Near a crossing of the
plane of a single reciprocal vector $b$, projecting onto the two crossing
momentum states gives diabatic splitting $\Delta E = |q+b-Q|^2 - |q-Q|^2$ with
slew rate $|d\,\Delta E/dt| = 2\, |b \cdot \dot Q|$ and gap $A_b$ (the
plane's own cosine amplitude, i.e.\ twice its complex Fourier coefficient
--- no other cosine contributes at first order). The
Landau-Zener loss per crossing is therefore
\begin{equation}
P_{\text{tunnel}} = \exp\!\left(-\frac{\pi A_b^2}{4\, |b\cdot \dot Q|}\right),
\qquad b\cdot\dot Q \neq 0,
\label{projlaw}
\end{equation}
which depends only on the component of the sweep \emph{normal to the plane}:
an oblique shuttle pays only for the normal velocity. For $d{=}1$,
$b = k$, $\dot Q = g/2$ recovers the manuscript's
$\exp(-\pi A^2/2kg)$.

\emph{Verification, separable case.} Cubic lattices in $d = 2, 3$
($b_i = 2\hat e_i$, amplitudes $A \in [0.2, 0.45]$, drives
$g \in [0.12, 0.24]$), sweep tilted by $\theta \in [0^\circ, 30^\circ]$ from
a plane normal, integration on $(2n{+}1)^d$ momentum grids: all 30 points
match eq.~(\ref{projlaw}) to $3\times10^{-2}$ relative or better, confirming
the $\cos\theta$ collapse. These lattices are separable ($H = \sum_i H_i$), so
this is a consistency check of the framework, not a test of the law beyond
1D physics (endpoint effects contribute at the $10^{-2}$ level, measured by
margin extrapolation).

\emph{Verification, nonseparable case.} A sharper test is a minimal
2D lattice --- three lasers, cosines on $b_1$, $b_2$ \emph{and} the cross
vector $b_1 - b_2$ --- which is not a Kronecker sum. There
eq.~(\ref{projlaw}) survives only as the leading-order description: measured
single-plane losses exceed the bare-$A$ prediction systematically, by
$+2$--$8\%$ relative at $A = 0.25$ and $+8$--$32\%$ at $A = 0.4$ (worst at
large tilt). The cross term dresses the crossing --- indirect paths through
the neighboring momentum state modify the effective gap --- so in real
minimal-laser lattices the projection law is a design estimate good to tens
of percent at moderate depth, not an exact statement.

\section{Simultaneous crossings at the zone corner}

Sweeping along the diagonal $\theta = 45^\circ$ crosses two Bragg planes at
the same instant. This is a \emph{four}-state event, not three: at the zone
corner $|q\rangle$, $|q{+}b_1\rangle$, $|q{+}b_2\rangle$ \emph{and}
$|q{+}b_1{+}b_2\rangle$ are simultaneously degenerate, and the fourth state
cannot be dropped --- it is what the ground band becomes past the corner. A
three-state truncation predicts a collective $\sqrt2$-enhanced coupling and
hence probability $P_1^2$ of remaining in the initial diabatic momentum
state (that model has no post-corner ground-band channel at all, which is
precisely its defect); integrating the reduced four-state model instead
gives carried population $\approx (1-P_1)^2$, and the full lattice
simulation agrees:
across 6 depth/drive combinations (losses $0.29$--$0.97$) the measured loss
matches the factorized prediction $1-(1-P_1)^2$ to $\lesssim 10^{-2}$
(median $\sim 1\times10^{-3}$, down to $\sim 10^{-4}$ at the highest
losses). The mechanism is separability, not a generic
multilevel theorem: the cubic potential splits as $H = H_x + H_y + H_z$, so
the diagonal sweep is exactly two independent 1D Landau-Zener events. For
\emph{nonseparable} lattices the factorization breaks, and we measured the
breakdown: in the minimal 3-laser 2D lattice (cross term
$\cos((b_1{-}b_2)\cdot x)$ coupling the degenerate intermediate states), the
diagonal-corner loss exceeds the factorized prediction by $+0.01$ to $+0.18$
absolute (largest at $A = 0.4$, $g = 0.15$), while the separable reference
at identical parameters agrees within $0.004$. Corners of real
minimal-laser lattices are anomalously lossy and deserve their own
treatment.

\section{Loss-reducing shuttle ramps}

Fix total transported momentum $Nk$ and total time $T$. Equation
(\ref{projlaw}) says loss is paid only inside windows of width
$w \sim \mathcal{O}(A/k)$ around the $N$ crossings, which suggests a bang-bang-like profile: sweep slowly
($\dot Q = v_s$) through each window, sprint ($\dot Q = v_f$) between them.
As $v_f/v_s \to \infty$ the budget concentrates in the windows,
$v_s \to N w / T$, and the per-crossing exponent improves over the linear
ramp by the factor $k/w > 1$. The window cannot be taken arbitrarily narrow:
the LZ transition region has its own sweep-rate-dependent width, and the
formula requires the crossing to be traversed between asymptotic detunings,
so the improvement cannot be unbounded; where the true optimum of the exact
dynamics lies, and under what control constraints, is left open.

\emph{Measurement} ($A = 0.35$, $k = 2$, $N = 3$, $T = 254$ chosen so the
linear ramp loses 34.1\%; controls restricted to a two-parameter
tanh-smoothed bang-bang family with bounded $v_f$): the best profile in a
12-point scan ($v_s/v_f = 0.15$, $w = 0.5$) achieves $0.28\%$ measured loss,
a reduction by a factor of 122 at equal total time, with the trend in
$(v_s/v_f, w)$ as expected over the probed range. Two caveats. First, the
independent-LZ estimate at the measured minimum crossing speed
($v_s = 0.0072$) is $0.37\%$; the measured value sits slightly below it,
within finite-window corrections, so the improvement is essentially
incoherent physics (a factor of 90 is supported by the incoherent estimate
alone). Second, this is the best of a small in-family scan, not a
demonstrated global optimum: loss-reducing, not loss-optimal.
The linear-ramp measurement itself validates the compounding law
$1-(1-P)^N$ to three decimals, and the best-profile loss is stable under
doubling the time resolution together with a $1.5\times$ increase of the
momentum cutoff.

\section{Uniform-depth lattices from shaped illumination}

Lattice depth is proportional to local beam intensity, so a Gaussian-beam
lattice is shallowest exactly where the cloud ends: LZ loss, exponential in
$A^2$, then eats the cloud edge first. Shaping the lattice beams into a
flat-top (via the optimal-transport method of Swan et al., Opt.\ Express 33,
6290 (2025); reimplemented and cross-checked here) equalizes the depth. With
equal power delivered into the cloud region, a simulated $\sigma = 18\,$px
cloud in the reconstructed flat-top vs an equal-power Gaussian shows:
cloud-averaged 3-crossing survival $0.863$ vs $0.846$ (comparable), depth
coefficient of variation $5.7\%$ vs $25.4\%$, and 5th-percentile survival
$0.882$ vs $0.564$. For the flat-top the mean sits \emph{below} the 5th
percentile; this is not an error but the shape of the distribution --- the
small fraction of the cloud beyond the plateau edge is lost almost entirely
while everything inside survives alike, which is exactly the sense in which
uniform illumination protects the bulk rather than raising the average. It
should also reduce sensitivity to shot-to-shot cloud-position fluctuations,
though we did not simulate a fluctuating ensemble.

A remark on the shaper itself. The baseline worth measuring against is not a
randomly initialized Gerchberg-Saxton loop but an analytically preconditioned
one: the standard package (\texttt{slmsuite}) seeds from a lens-and-blaze fit
to the target's first and second moments, which is much better than random.
Against all three on a $256\,$px aperture padded to $512$, after $50$
iterations, scored with that package's own winding-number measure (random
initialization averaged over five seeded draws): residual vortices in the
flat-top $300$ (random), $32$ (analytic), $0$ (OT-seeded); intensity ripple
$0.347$, $0.340$, $0.178$; efficiency $0.927$, $0.959$, $0.981$. The ripple
floor set by the target's own soft edges is $0.175$, so the OT seed lands
essentially at the best the target admits. Brenier's theorem motivates the
construction --- the exact transport map is the gradient of a convex potential
--- but the zero winding numbers are measured rather than derived: the discrete
solver only approximates that map, and a curl-free aperture phase does not by
itself force a vortex-free far field. The advantage is insensitive to what runs
afterwards. Across six Gerchberg--Saxton variants, two targets, and plain and
noise-region variants of each, the OT seed gives zero vortices in all $24$
configurations, while the analytic seed ranges from $0$ to $125$ depending on
the variant. On a ring target the analytic seed already reaches zero vortices
unaided and the two tie. At practical coarse-grid settings the Sinkhorn solve
costs $0.3\,$s, comparable to the iterations it precedes.

\section{The rank criterion for optical quasicrystals}

Swan's manuscript observes that non-minimal laser configurations are
``generically not a lattice'' and asks when they form quasicrystals. The
sharp criterion: let $M \subset \mathbb{R}^d$ be the $\mathbb{Z}$-module
generated by the potential's wavevectors,
$r = \operatorname{rank}_{\mathbb Z} M$, and
$e = \dim_{\mathbb R} \operatorname{span} M$. On its span, the potential is
periodic iff $r = e$ (the module is then a true reciprocal lattice); $r > e$
gives a \emph{quasiperiodic potential} with an $r$-dimensional superspace
description (we use ``quasicrystal'' informally for the potential; properties
of the resulting atomic phase --- long-range order, pure-point diffraction,
stability --- are separate questions not addressed here).
(Comparing $r$ to $d$ alone fails: $\cos x + \cos\sqrt2\,x$ in the plane has
$r = 2 = d$ but $e = 1$ --- quasiperiodic along $x$, trivial in $y$.) The rank
statements are algebraic; the four configurations below illustrate them
numerically, using a bounded-denominator rational-dependence heuristic
(floating-point rank alone cannot certify integer relations, and exact
certification for general generator sets would need lattice-reduction
methods; the four configurations here are standard ones whose module ranks
are known analytically): a minimal
3-laser 2D lattice (rank 2 = $d$:
periodic), the pentagonal 5-laser configuration (rank 4 $> 2$: quasicrystal,
tenfold diffraction), a 4-laser configuration with the redundant beam chosen
on the integer lattice per Swan's countable-family construction (rank 2:
periodic, as he predicted), and the icosahedral configuration in 3D, realized as six standing waves
along the icosahedral axes ($r = 6 > e = 3$: the icosahedral quasicrystal;
six independent standing waves require frequency offsets \emph{between
different pairs} --- each pair internally monochromatic --- so that
cross-pair interference terms time-average away, per the manuscript's
multi-frequency remark). His open question --- which
space groups are realizable --- remains open and is a natural next target;
the Bravais-class half is settled by his own construction.

\section*{Reproducibility}

\texttt{lattice3D.py}, \texttt{run\_3d.py}, \texttt{ramp\_optimizer.py},
\texttt{cloud\_survival.py}, \texttt{phase\_diversity.py},
\texttt{quasicrystal.py}; results JSON alongside. The phase-diversity solver
(cf.\ the empty \texttt{PhaseDiversity} repository) recovers a simulated
aberrated input beam from 6 defocus probes with overlap
$|\langle E_{\rm rec}|E_{\rm true}\rangle|^2 = 0.9997$ under Poisson shot
noise at $10^4$ peak counts per pixel (0.8\% amplitude error;
intensity-weighted residual phase RMS 12~mrad, or 28~mrad unweighted over
the region with amplitude above 5\% of peak --- the dim edge dominates the
unweighted figure).

\end{document}
